\documentclass[12pt,a4paper]{article}
\usepackage[utf8]{inputenc}
\usepackage[french]{babel}
\usepackage[T1]{fontenc}
\usepackage{geometry}
\geometry{margin=2.5cm}

\title{Conclusion Générale du Projet}
\date{}

\begin{document}

\maketitle

\section{Synthèse du projet}
Le projet a permis de concevoir un système complet de gestion du cahier de texte. Nous avons atteint les objectifs de dématérialisation et de sécurisation des données pédagogiques, facilitant ainsi la communication entre tous les acteurs de l'établissement.

\section{Récapitulatif de l’architecture}
La solution repose sur une \textbf{architecture 3-tiers} moderne avec un Frontend React/Angular, un Backend Spring Boot/Node.js et une base de données PostgreSQL assurant l'intégrité des données.

\section{Analyse de qualité}
Le système répond aux exigences de \textbf{sécurité} (AuthService), de \textbf{traçabilité} (AuditService) et de \textbf{disponibilité}. La séparation des services permet une maintenance aisée.

\section{Recommandations pour le développement}
Nous suggérons un déploiement en 5 phases :
\begin{itemize}
    \item \textbf{Phase 1 (S1-4) :} Fondations (Infrastructure et AuthService).
    \item \textbf{Phase 2 (S5-10) :} Fonctionnalités Core (Sessions et Classes).
    \item \textbf{Phase 3 (S11-14) :} Validation Admin et Audit.
    \item \textbf{Phase 4 (S15-18) :} Rapports PDF et Notifications.
    \item \textbf{Phase 5 (S19-20) :} Mise en production et formation des utilisateurs.
\end{itemize}

\section{Extensions futures}
Les évolutions prévues incluent le développement d'une application mobile pour les appels en classe et l'intégration de modules d'intelligence artificielle pour l'analyse prédictive du décrochage scolaire.

\end{document}
